\cotefile{37}
\scanpage{58}
Pour Exp $\overline{\underline{\mathrm{II}}}$, \inserted{2.2.6} signaler

SGA 2 \quad $\overline{\underline{\mathrm{V}}}$ 3.5.

p. 10.
% archive sheet number: 57

\scanpage{59}
Posons
\[
 \bar i\,\bar j=\tau
\]

Pour $\forall\ E\in\operatorname{Ob}C$, \unsure{nous allons} définir une application
\[
 \text{\textcircled{*}}\qquad
 \tau(E)\times E\xrightarrow{\varphi_E}E
\]
qui fera de $E$ un objet formellement principal sous l'objet-groupe $\tau(E)$.
\deleted{Je \unsure{prends}} \inserted{\unsure{Pour ceci|Pour voir}} si \ldots

\[
 \deleted{\text{\unsure{on|ou}}\quad\unsure{\bar{\imath}|=}\quad\unsure{\cdot|i}}\quad
 \cdot\longrightarrow J\longrightarrow E\xrightarrow{{\unsure{p|\rho|\varphi}}_E}B
\]
\unsure{on aura|on a} $\unsure{\tau(E)|\tau(F)}=D_{\unsure{B|E}}(\unsure{J|\mathcal{J}})$,
% archive sheet number: 58

\scanpage{60}
\begin{center}
INSTITUT\\
DES HAUTES ÉTUDES SCIENTIFIQUES\\[-.4ex]
\rule{1.4cm}{.4pt}
\end{center}

SGA 6 $\overline{\underline{\mathrm{VII}}}$ \qquad \unsure{manque|marquer} la formule \unsure{de|des} self-\unsure{intersection|intersections} pour $\mathrm{Gr}(X),\mathrm{Gr}(Y)$

(après \deleted{\unsure{6|2|4}}4.\deleted{\unsure{5|6}}\unsure{6|5})

Peut-être fin de l'Exp ?
% archive sheet number: 53

\scanpage{61}
\begin{center}
INSTITUT\\
DES HAUTES ÉTUDES SCIENTIFIQUES\\[-.4ex]
\rule{1.4cm}{.4pt}
\end{center}

SGA 6

Structure additive explicite
de l'anneau $\unsure{\mathrm{Gr}}(\widetilde X)\otimes_{\unsure{Z|\mathbb Z}}\unsure{Q|\mathbb Q}$ d'un
schéma éclaté..---
% archive sheet number: 60

\scanpage{62}
\begin{center}
INSTITUT\\
DES HAUTES ÉTUDES SCIENTIFIQUES\\[-.4ex]
\rule{1.4cm}{.4pt}
\end{center}

\uline{Exp}. X \qquad Tenir compte conventions \unsure{de|des} Illusie (faire relire par Berthelot)
% archive sheet number: 61

\scanpage{63}
\begin{center}
INSTITUT\\
DES HAUTES ÉTUDES SCIENTIFIQUES\\[-.4ex]
\rule{1.4cm}{.4pt}
\end{center}

SGA 6 $\overline{\underline{\mathrm{III}}}$ 1, p. 15, signaler en footnote
\margin{\deleted{\unsure{15|14}}}qu'un morphisme plat loc. de pres. finie \unsure{est|et}
parfait, et à \unsure{fortiori|fortiori.} pseudo-cohérent.
% archive sheet number: 62

\scanpage{64}
\begin{center}
INSTITUT\\
DES HAUTES ÉTUDES SCIENTIFIQUES\\[-.3ex]
\rule{1.3cm}{.4pt}
\end{center}

{\reversemarginpar\margin{$\unsure{\alpha|\times}$ AW}}
\begin{tabular}{@{}r@{\quad}l@{}}
p. 43 & $(\overline{\mathrm{I}\underline{\mathrm{X}}}\quad\ldots\quad)$\\
p 47 & $(\overline{\underline{\mathrm{IV}}}\qquad)$\\
p. 51 & $(\overline{\underline{\mathrm{IV}}}\qquad)$\\
p. 52 & $(\overline{\underline{\mathrm{IV}}}\qquad)$\\
p. 56 & $(\overline{\underline{\mathrm{IX}}}\qquad)$\\
p. 59 & $(\overline{\underline{\mathrm{IV}}}\qquad)$\\
p. 70 & $(\mathrm{IX}\qquad)$\\
p. \unsure{74|71} & $(\overline{\underline{\mathrm{IX}}}\qquad)$\\
p. 75 & $(\text{\unsure{id|iv}}\qquad)$
\end{tabular}
% archive sheet number: 63

\scanpage{73}
\begin{flushright}
SGA 7
\end{flushright}
% archive sheet number: 72

\scanpage{78}
SGA 7 $\underline{\overline{\mathrm{VI}}}$ App

après 1.2 harmoniser notation $L\Omega_{\unsure{X_{\varphi}|X/\varphi|X/S}}$ pour compl. cot. rel.

mettre bibliographie commune \unsure{:|à} l'exposé et l'App.
% archive sheet number: 77

\scanpage{79}
\begin{center}
INSTITUT\\
DES HAUTES ÉTUDES SCIENTIFIQUES\\[-.3ex]
\rule{1.3cm}{.4pt}
\end{center}

\deleted{SGA 7 $\underline{\overline{\mathrm{VIII}}}$\\
Faux schémas abéliens (ou \unsure{Motifs|motifs} mixtes\\
effectifs \unsure{de|des} poids 1)\\
Éléments de N.S. et correspondances symétriques.}

\noindent\rule{.62\linewidth}{.4pt}

$\underline{\overline{\mathrm{I}}}$ \quad Le groupe d'inertie $I$ a un\\
\hspace*{2em}$p$-ss-gpe de Sylow infini

Cf Serre \unsure{Corps locaux|Corps claux} \unsure{Coh galois|Cohom galois}, p.144, lemme 4.

\noindent\rule{.88\linewidth}{.4pt}

$\underline{\overline{\unsure{\mathrm{VI}|\mathrm{VII}}}}$ \quad Traiter aussi le ``local invariant\\
\hspace*{2em}cycle problem'' pour \deleted{\unsure{des} \illegible{1}}\inserted{$H^{1}$}

Sign. \unsure{la|le} conjecture de Griffiths en dim. quelc.

\noindent\rule{.96\linewidth}{.4pt}
% archive sheet number: 78

\scanpage{80}
\begin{center}
INSTITUT\\
DES HAUTES ÉTUDES SCIENTIFIQUES\\[-.3ex]
\rule{1.3cm}{.4pt}
\end{center}

\noindent
\begin{minipage}[t]{.55\linewidth}
\begin{tabular}[t]{@{}l@{\;}r@{\quad}l@{\quad}l@{}}
SGA 7 $\overline{\underline{\mathrm{IX}}}$ & (iii) & l. $-5$ & \\
 & 1 & l. 5 & \\
 & 33 & l. 10 & \\
 & 35 & l. 5 & \\
 & 37 & l. 8 & \\
 & 69 & l. $-3$ & \\
 & 70 & l. 2 & \\
 & 73 & l. 12 & \\
 & 87 & l. 6 & livre de Gabriel-Demazure\\
 & 124 & l. $-8$ & \\
 & 128 & l. 4, l. $-9$ & \\
 & 134 & l. 1 & \\
 & 146 & l. 5 & \\
 & 150 & l. $-3$ & \\
 & 157 & l. 5 & \\
 & 1\unsure{\deleted{7}5|\deleted{5}5|5}9 & l. 11, l. $-7$ & \\
 & 169 & l. 4 & deux fois
\end{tabular}
\end{minipage}\hfill
\begin{minipage}[t]{.42\linewidth}
\begin{tabular}[t]{@{}c@{\quad}l@{}}
\uline{Bibl}. & \unsure{ajouter} [22 bis]\\
{[13]} & \\
{[24]} &
\end{tabular}
\end{minipage}
% archive sheet number: 79

\scanpage{81}
\begin{center}
INSTITUT\\
DES HAUTES ÉTUDES SCIENTIFIQUES\\[-.3ex]
\rule{1.3cm}{.4pt}
\end{center}

SGA 7 $\underline{\overline{\mathrm{XVII}}}$

\begin{tabular}{@{}ll@{}}
p. 2 & 1.\\
p. 7 & l. $-7$\\
p. 8 & l. 17\\
p. 24 & l. 8
\end{tabular}
% archive sheet number: 80

\scanpage{82}
\begin{center}
INSTITUT\\
DES HAUTES ÉTUDES SCIENTIFIQUES\\[-.3ex]
\rule{1.3cm}{.4pt}
\end{center}

Appendice à $\underline{\overline{\mathrm{XXI}}}$ de Katz : torsion en coh. cristalline

$\underline{\overline{\mathrm{XXII}}}$ \quad Référence à la démonstration de Deligne du th. de Congruences de Katz (\unsure{généralisée|génér. lisse})

{\reversemarginpar\margin{$\underline{\overline{\unsure{\mathrm{VI}|\mathrm{XI}}}}$\\ou $\underline{\overline{\mathrm{XIII}}}$\unsure{?|'}}}Vérifier qu'on y trouve le fait qu'en déformant une singularité quadratique ordinaire, \unsure{---|on} \unsure{donne} au plus autant (y inclus cas de car. rés. 2 et dim relative paire). En d'autres termes, l'ens des pts où un morphisme plat \deleted{de}\inserted{loc.} prés. finie est lisse ou a une singularité quadratique ordinaire est tjrs. \uline{ouvert}.
% archive sheet number: 81

\scanpage{83}
SGA 7 $\overline{\underline{\mathrm{XVIII}}}$ \qquad p. 10 l. 2 réf à \deleted{\unsure{SGA 7|SGA 2|SGA|Exp} $\overline{\underline{\unsure{\mathrm{XVIII}|\mathrm{XIII}}}}$ Exp} SGA 5 $\overline{\underline{\mathrm{IV}}}$

\begin{tabular}{@{}l@{\quad}l@{\quad}l@{}}
p. 17 & l. 9 & réf \unsure{Exp|l'Exp} $\overline{\underline{\mathrm{XIV}}}$\\
p 25 & l. $-2$ & ---\quad SGA 5 $\overline{\underline{\mathrm{IV}}}$\\
p. 34 & l 4 & ---\quad $\overline{\underline{\mathrm{XVI}}}$\\
p. 44 & l, 1 & \\
p 64 & l. $-7$, l. $-2$ & --\quad $\overline{\underline{\mathrm{XVI}}}$\\
p. 66 & l. 6 & ---\quad $\overline{\underline{\mathrm{XIII}}}$\\
 & l. $-5$ & id ?\\
p. 68 & l. 4 & $\overline{\underline{\mathrm{XVI}}}$\\
 & l. $-6$ & $\overline{\underline{\mathrm{XII}}}$
\end{tabular}
% archive sheet number: 82

\scanpage{84}
\begin{flushright}
\uline{Réédition EGA I}\\
\deleted{\unsure{et} \illegible{1}}
\end{flushright}
% archive sheet number: 83

\scanpage{95}
\begin{center}
INSTITUT\\
DES HAUTES ÉTUDES SCIENTIFIQUES\\[-.3ex]
\rule{1.3cm}{.4pt}
\end{center}

10.12. Stabilité \unsure{par} $\underset{\longleftarrow}{\lim}$ (ou non ?)
% archive sheet number: 94

\scanpage{99}
\begin{flushright}
EGA $\overline{\underline{\mathrm{II}}}$
\end{flushright}
% archive sheet number: 98

\scanpage{100}
\begin{center}
INSTITUT\\
DES HAUTES ÉTUDES SCIENTIFIQUES\\[-.3ex]
\rule{1.3cm}{.4pt}
\end{center}

\textcircled{a} \quad \unsure{transférer} EGA $\underline{\overline{\mathrm{IV}}}$ 4.8
\[
\left\{\begin{aligned}
&\text{Également suite des}\\
&\text{morphismes affines \unsure{des|de}}\\
&\text{\unsure{préschémas|préschèmes} formels.}
\end{aligned}\right.
\]

{\reversemarginpar\margin{!}}Appendice \unsure{par|pour} Murre aux schémas formels ?

\textcircled{b} \quad Cf EGA $\underline{\overline{\mathrm{II}}}$ 6.5.10 ajouter formules de la norme

6.5.11 ajouter \unsure{formules|formule} pour le cas où $X'\to X$ est gén. \unsure{décomposé} (du moins \unsure{géom!|géomt.})

6.5.\unsure{9|1} renvoi à EGA $\underline{\overline{\mathrm{IV}}}$ 21..

\textcircled{c} \quad Remonter EGA $\underline{\overline{\mathrm{IV}}}$ 19.4.11

(revu et augmenté)

\textcircled{d} \quad 7.4.\unsure{2|3}. Admettre courbe vide ---
% archive sheet number: 99
